\chapter{Project summary}
\label{app:project-summary}

\section*{Title}

\textit{A Hard Negative-Story Threshold: Training-Period Association, Temporal Non-Replication and Economic Limits}

\section*{Problem}

Financial-news systems often average every sentiment score for a firm on a given
day. Counting the share of stories classified as negative gives a different
measure. The project tests whether that measure is related to returns after the
main controls. It also asks when the relationship appears and whether it remains
useful in later data and after trading costs.

\section*{Objectives}

\begin{enumerate}
  \item Build a reproducible firm-day dataset joining financial news, sentiment scores and assigned-session open-to-open returns.
  \item Compare several return-blind rules for aggregating multiple same-day stories.
  \item Test whether the selected measure is still related to returns after controlling for average sentiment, news volume and recent prices.
  \item Separate the assigned return into intraday and overnight parts and state the timestamp limitation.
  \item Evaluate selection, uncertainty, multiple testing, transaction costs and downside-risk performance.
  \item Test the measure on a separate Reuters/LSEG company set without combining the two data sources.
\end{enumerate}

\section*{Methods}

FinBERT, a language model for financial sentiment, scores each story. Daily rank
correlations and rank regressions test the relationship. The statistical methods
allow for links between nearby sessions and for the fact that several measures
are tested. The economic tests include turnover, trading costs, a volatility-based
risk rule, controls with similar market exposure and checks that shift the risk
schedule to other dates.

\section*{Outcome}

Negative-story share has a small exploratory relationship with returns in the
training period. The relationship appears during the assigned open-to-next-open
window and in an earlier return window. Because exact story availability times
are missing, the timing evidence does not support a forecasting claim. The
unchanged specification produces an opposite-signed estimate during 2020--2023.
This suggests instability but does not explain its cause. Daily trading does not
cover costs. In FNSPID, the same-session news risk rule has a worse overall maximum
drawdown than its volatility benchmark. It also does not show that the news was
available before the return began. The LSEG estimates are too uncertain, and the
tuned LSEG risk and prompt rules fail their economic tests.
